% Part: computability
% Chapter: machines-computations
% Section: configuration

\documentclass[../../../include/open-logic-section]{subfiles}

\begin{document}

\olfileid{cmp}{tur}{con}
\olsection{સંરૂપણો અને સંગણનાઓ}

\begin{explain}
અગાઉ આપણે સમ મશીનનાં સંરૂપણોનો ક્રમ અનુસર્યો હતો તે યાદ
કરો. ટેપ, વાંચન-લેખન હેડ અને ટ્યુરિંગ મશીનના પ્રોગ્રામથી
બનેલું કલ્પિત ઉપકરણ ખરેખર ટ્યુરિંગ મશીનની સંગણના શું છે
તેની સહજ ચિત્રાત્મક સમજ આપવાની રીત છે. ઔપચારિક રીતે,
આપેલા ઇનપુટ પર ટ્યુરિંગ મશીનની સંગણનાને
\emph{સંરૂપણો}ના અનુક્રમ તરીકે વ્યાખ્યાયિત કરી શકાય છે.
સંરૂપણ પોતે પ્રતીકોનો અનુક્રમ (સંગણનાના કોઈ સમયે ટેપ
પર શું લખેલું છે તે દર્શાવતો), વાંચન-લેખન હેડનું સ્થાન
જણાવતી સંખ્યા અને એક અવસ્થા છે. આ બધાની મદદથી આપણે
આપેલા ઇનપુટ પર ટ્યુરિંગ મશીન~$M$ શું ગણે છે તે
વ્યાખ્યાયિત કરી શકીએ છીએ.
\end{explain}

\begin{defn}[સંરૂપણ]
ટ્યુરિંગ મશીન $M = \tuple{Q, \Sigma, q_0,
\delta}$નું \emph{સંરૂપણ} એ $\tuple{C, m, q}$ ત્રિક છે, જ્યાં
\begin{enumerate}
\item $C \in \Sigma^*$ એ $\Sigma$નાં પ્રતીકોનો સાન્ત અનુક્રમ છે,
\item $m \in \Nat$ એ $< \len{C}$ એવી સંખ્યા છે, અને
\item $q \in Q$ છે.
\end{enumerate}
સહજ રીતે, અનુક્રમ~$C$ ટેપની સામગ્રી છે (સૌથી ડાબા
ખાનાથી માંડીને છેલ્લાં ખાલી ન હોય અથવા અગાઉ મુલાકાત
લીધેલાં ખાનાં સુધીનાં પ્રતીકો); $m$ વાંચન-લેખન હેડ જે
ખાનું વાંચે છે તેનો ક્રમાંક છે (સૌથી ડાબા ખાનાનો ક્રમાંક
$0$ રાખીને), અને $q$ મશીનની વર્તમાન અવસ્થા છે.
\end{defn}

\begin{explain}
ટ્યુરિંગ મશીનનું સંભવિત ઇનપુટ પ્રતીકોનો અનુક્રમ હોય છે,
સામાન્ય રીતે કોઈ સ્વરૂપમાં એક સંખ્યાનું સંકેતીકરણ કરતો
અનુક્રમ. ટ્યુરિંગ મશીનનું પ્રારંભિક સંરૂપણ એ સંરૂપણ છે જેમાં
મશીનને તે ઇનપુટ પર કામ કરવા શરૂ કરીએ છીએ: ટેપ પર પહેલાં
ટેપના છેડાનું ચિહ્ન હોય છે અને તેની તરત જમણેનાં ખાનાંમાં
ઇનપુટ લખેલું હોય છે; વાંચન-લેખન હેડ ઇનપુટનું સૌથી ડાબું
ખાનું (એટલે છેડાના ચિહ્નની જમણેનું ખાનું) વાંચે છે; અને
ઉપકરણ નિર્દિષ્ટ પ્રારંભિક અવસ્થા~$q_0$માં હોય છે.
\end{explain}

\begin{defn}[પ્રારંભિક સંરૂપણ]
ઇનપુટ $I \in \Sigma^*$ માટે $M$નું \emph{પ્રારંભિક સંરૂપણ}
આ છે:
\[
\tuple{\TMendtape \frown I, 1, q_0}.
\]
\end{defn}

\sourcecorrection{OLTM-001}{સ્થિર અંગ્રેજી મૂળનો પરિચય
શરૂઆતમાં હેડ ટેપનું સૌથી ડાબું ખાનું વાંચે છે એમ કહે છે;
અહીંની ઔપચારિક વ્યાખ્યા અને $m=1$નું સૂત્ર હેડને
ઇનપુટના પ્રથમ ખાનામાં, છેડાના ચિહ્નની જમણે, મૂકે છે.
આ વિભાગમાં પ્રારંભિક સંરૂપણ માટે ઔપચારિક સૂત્ર અનુસર્યું છે.}

\begin{explain}
ચિહ્ન~$\frown$ \emph{અનુક્રમોનું જોડાણ} દર્શાવે છે---ઇનપુટ
શબ્દમાળા ડાબા છેડાના ચિહ્નની તરત જમણે શરૂ થાય છે.
\end{explain}

\sourcecorrection{OLTM-002}{સ્થિર અંગ્રેજી મૂળ આ વાક્યમાં
ઇનપુટ છેડાના ચિહ્નની તરત ``ડાબે'' શરૂ થાય છે એમ કહે છે.
પરંતુ બાજુનું $\TMendtape \frown I$ સૂત્ર અને અગાઉનું
ટેપનું વર્ણન ઇનપુટને તેની જમણે મૂકે છે; ગુજરાતી ગદ્યમાં
તે મુજબ દિશા સ્પષ્ટ કરી છે.}

\begin{defn}
અપણે કહીએ છીએ કે (મશીન~$M$ મુજબ) સંરૂપણ
$\tuple{C, m, q}$ \emph{એક પગલામાં સંરૂપણ
$\tuple{C', m', q'}$ આપે છે}, ત્યારે અને ત્યારે જ જ્યારે
\begin{enumerate}
\item $C$નું $m$-મું પ્રતીક $\sigma$ છે,
\item $M$નો સૂચનાઓનો ગણ $\delta(q, \sigma) =
  \tuple{q', \sigma', D}$ નિર્દિષ્ટ કરે છે,
\item $C'$નું $m$-મું પ્રતીક $\sigma'$ છે, અને
\item
\begin{enumerate}
\item $D = L$ હોય, અને એ કિસ્સામાં $m>0$ હોય તો
  $m' = m - 1$, નહીંતર $m'=0$; અથવા
\item $D = R$ હોય અને $m' = m + 1$ હોય, અથવા
\item $D = N$ હોય અને $m' = m$ હોય,
\end{enumerate}
\item $m' = \len{C}$ હોય તો $\len{C'} = \len{C} + 1$ છે
  અને $C'$નું $m'$-મું પ્રતીક~$\TMblank$ છે. નહીંતર
  $\len{C'}=\len{C}$ છે.
\item $i < \len{C}$ અને $i \neq m$ એવા દરેક $i$ માટે
  $C'(i) = C(i)$ છે,
\end{enumerate}
\end{defn}

\begin{defn}\ollabel{defn:run-output}
\emph{ઇનપુટ~$I$ પર $M$નું ચલન} એ $M$નાં સંરૂપણોનો
અનુક્રમ $C_i$ છે, જેમાં $C_0$ એ ઇનપુટ~$I$ માટે
$M$નું પ્રારંભિક સંરૂપણ છે અને દરેક $C_i$ એક પગલામાં
$C_{i+1}$ આપે છે.

જો $C_k = \tuple{C, m, q}$ હોય, $C$નું $m$-મું
પ્રતીક~$\sigma$ હોય અને $\delta(q, \sigma)$ અવ્યાખ્યાયિત
હોય, તો આપણે કહીએ છીએ કે $M$ \emph{ઇનપુટ $I$ પર
$k$ પગલાં પછી થંભે છે}. એ કિસ્સામાં, ઇનપુટ~$I$ માટે
$M$નું \emph{આઉટપુટ} $O$ છે, જ્યાં $O$ એ
$\TMblank$માં પૂરું ન થતું એવું પ્રતીકોનું શબ્દમાળા છે કે
કોઈ~$j \in \Nat$ માટે
$C = \TMendtape \concat O \concat \TMblank^j$.
($\TMblank^j$ એ $j$ $\TMblank$નો અનુક્રમ છે.)
\end{defn}

\begin{explain}
આ વ્યાખ્યા મુજબ $M$નું આઉટપુટ~$O$ હંમેશાં
$\TMblank$ સિવાયના પ્રતીકમાં પૂરું થાય છે; અથવા
સમય~$k$ પર આખી ટેપ $\TMblank$થી ભરેલી હોય
(સૌથી ડાબા $\TMendtape$ સિવાય), તો $O$ ખાલી
શબ્દમાળા છે.
\end{explain}

\end{document}
